
\subsection{Scope of the generated Documentation}\label{sec:GenerationScope}
The scope of the generated documentation depends on the type of product you have to deliver (refer to \tqfullvref{sec:TypesOfProducts}) and the audience for that documentation.

The table below shows the command line options that determine the volume of the generated documentation for the product types and audiences.

\renewcommand{\cellalign}{tl}
\LTXtable{\linewidth}{comments/Scope.tbl.tex}

For a function library\index{product type!function library}, the \bdq{users} are external developers that integrate that library into their product. That is slightly different for a feature library\index{product type!feature library}: the users will just add the library, while external developers dig deeper (admittedly, this distinction is rather arbitrary).

For tools\index{product type!tool}, extensions\index{product type!extension}, standalone\index{product type!standalone application} and server-based applications\index{product type!server-based application}, no Javadoc\index{javadoc} documentation will be generated for the public. Only when a standalone application\index{product type!standalone application} allows customisation, it may make sense to generate a Javadoc\index{javadoc} documentation for the public API that can be used by external developers for the customisations.

\subsubsection{The Command Line Options}
The meaning of the command line options is described below:

\paragraph{\texttt{-author}} If set, the value of the \nameref{sec:TagAuthor} tag will be included into the generated documentation. It might be omitted for documentation meant for the public.

\paragraph{\texttt{-{}-date}} Sets the date when the documentation was generated. If omitted and \bsq{\texttt{-notimestamp}} was not set either, the current data and time is used. This is sometimes not wanted, as it makes the build non-reproducible.

\paragraph{\texttt{-{}-expand-requires}} Instructs the Javadoc\index{javadoc} tool to expand the set of modules to be documented. If not set, only the modules given explicitly on the command lin are documents. With the argument \bsq{\texttt{transitive}}, all the required transitive dependencies of those modules are included, and the argument \bsq{\texttt{all}} includes all dependencies. 

\paragraph{\texttt{-linksource}}\label{par:linksource} Adds an HTML\index{HTML} version of each source file to the documentation and links the respective program elements to it. As it reveals all the details, it is usually used only for internal documentations. The option \bsq{\texttt{-sourcetab}} determines the number of blank that should be used to replace tab characters in the source.

\paragraph{-nodeprecated} Prevents the generation of the documentation for any deprecated element. It might by useful for function\index{product type!function library} and feature libraries\index{product type!feature library} when you want to prevent that deprecated APIs are still used. 

\paragraph{-notimestamp} Supresses the time stamp, which is hidden in an HTML\index{HTML} comment in the generated documentation.

\paragraph{\texttt{-private}, \texttt{-protected}, \texttt{-public}} These three options\footnote{There is also a fourth one, \mbox{\bsq{\texttt{-package}}}, that is usually not used.} controls what is shown in the generated documentation: if \mbox{\bsq{\texttt{-private}}} is set, all classes and members are part of the documentation, \mbox{\bsq{\texttt{-protected}}} means that only protected and public classes and members will be included, and \mbox{\bsq{\texttt{-public}}} is for the public classes and members only. Obviously, these options are mutually exclusive.

\paragraph{\texttt{-{}-show-members}} Specifies which members (fields, methods, constructors) are documented. If specified, it overrides \mbox{\bsq{\texttt{-private}}}, \mbox{\bsq{\texttt{-protected}}} and \mbox{\bsq{\texttt{-public}}} in regards of the members.

\paragraph{\texttt{-{}-show-module-contents}} Specifies the documentation granularity of module declarations; with the value \bsq{\texttt{api}} (the default) only the exported packages of a module are displayed in the module overview. Internal, non-exported packages are hidden unless they are included via other options. This corresponds to the view of an external user who sees the module only through its public API.

The value \bsq{\texttt{all}} displays all of the module's packages in the documentation –including non-exported (internal) ones – along with additional module information such as \texttt{requires} clauses for all dependencies (including transitive and optional ones), \texttt{uses} and \texttt{provides} clauses for services, and so on. This is the view required for sustaining and maintenance, or when working on the module implementation itself.

\paragraph{\texttt{-{}-show-packages}} Specifies which packages show up in the module documentation; possible values are \bsq{\texttt{exported}} for only the exported packages, or \bsq{\texttt{all}}.

\paragraph{\texttt{-{}-show-types}} Specifies which types (classes, interfaces, records, enums, annotations) are included into the generated documentation. If specified, it overrides \mbox{\bsq{\texttt{-private}}}, \mbox{\bsq{\texttt{-protected}}} and \mbox{\bsq{\texttt{-public}}} in regards of the types.

\paragraph{\texttt{-sourcetab}} see \bsq{\nameref{par:linksource}} \hyperref[par:linksource]{above}.

\paragraph{\texttt{-subpackages}} The packages listed here and recursively their subpackages will be included into the generated documentation. The option \bsq{\texttt{-exclude}} allows to eliminate packages from that list. In both lists, the colon (\bdq{:}) is used to separate the packages.

The internal documentation for testers, sustaining and maintenance staff it is sufficient to just give the name of the root package here, because it needs to contain all packages.

For the public documentation you must decide individually which packages should be revealed.

%==============================================================================
% End of file!!
%==============================================================================

